% !TeX root = ../../Book5.tex
% 纲要标题：### 21.2 代数群的坐标环
% 纲要位置：工作记录/计划纲要.md，第 4397 行。

\section{代数群的坐标环}
\label{b5:sec:2102}


% 纲要范围：
%
% * 仿射代数群
% * 群乘法与 Hopf 结构
% * 一般线性群及可对角化群
% * 特殊线性群到射影一般线性群的中心商与表示下降
%

代数群的乘法也可由函数上的拉回来描述。采用坐标环而非只看一个基域上的点，有两个好处：公式可以同时在不同交换代数上取值，正特征中的幂零信息也不会因为点集太小而消失。

\subsection{仿射代数群}
\label{b5:subsec:2102:01}

\begin{definition}\label{b5:def:affine-group}
设 \(k\) 为域，\(A\) 为有限生成的交换 \(k\) Hopf 代数，结构映射记为 \(\Delta,\varepsilon,S\)。对每个交换含幺 \(k\) 代数 \(R\)，记 \(\eta_R:k\to R\) 为单位映射，并令
\[
G(R)=\operatorname{Hom}_{k\text{-alg}}(A,R),
\]
对 \(a\in A\)，用 Sweedler 记号写 \(\Delta(a)=\sum a_{(1)}\otimes a_{(2)}\)。对 \(g,h\in G(R)\)，用余乘法、余单位及对极分别定义
\[
(gh)(a)=\sum g\bigl(a_{(1)}\bigr)h\bigl(a_{(2)}\bigr),\qquad
1_G=\eta_R\varepsilon,\qquad g^{-1}=g\circ S.
\]
这个群值函子称为一个有限型仿射群概形，记为 \(G=\operatorname{Spec}A\)，其坐标环记为 \(k[G]=A\)。在代数闭域上，若坐标环约化，也把相应的仿射簇及其群运算称为\term[fangshe-daishuqun]{仿射代数群}[affine algebraic group]。
\end{definition}
因为 \(A,R\) 都交换，两个代数同态的卷积仍为代数同态；Hopf 公理正好使以上运算结合、有单位与逆。一个 \(k\) 代数同态 \(R\to R'\) 与这些公式交换，所以这是对 \(R\) 自然的群，而非彼此不相干的一组点集。
\ProseParagraphBreak
两个这类群之间的态射，定义为反向的 Hopf 代数同态。特别地，点上的映射必须在所有 \(R\) 上由同一坐标公式给出；只给 \(G(k)\to H(k)\) 的一个抽象群同态，并不足以定义代数群态射。

\subsection{群乘法与 Hopf 结构}
\label{b5:subsec:2102:02}

坐标环方法与通常的正则映射方法相符。设 \(k\) 代数闭，\(G\) 为仿射簇，乘法、单位及求逆都是正则映射。则乘法拉回给出
\[
\Delta:k[G]\to k[G\times G]=k[G]\otimes k[G],
\]
单位与求逆分别给出 \(\varepsilon,S\)。正则函数在三元组 \((g,h,\ell)\) 上比较时，群乘法结合律就是余结合律，其他公理同理。反向从约化的有限型交换 Hopf 代数出发，正则函数关系恢复同样的运算。
\begin{proposition}\label{b5:prop:hopf-ideals-subgroups}
设 \(k\) 为域，\(A=k[G]\) 为有限生成的交换 \(k\) Hopf 代数，其结构映射为 \(\Delta,\varepsilon,S\)。若理想 \(I\subseteq A\) 满足
\[
\Delta(I)\subseteq I\otimes A+A\otimes I,\qquad
\varepsilon(I)=0,\qquad S(I)\subseteq I,
\]
则 \(A/I\) 自然成为交换 Hopf 代数，且
\(\operatorname{Spec}(A/I)\) 是 \(G\) 的闭子群概形。反过来，使商上的三种映射良定义所需的条件正是以上三式。
\end{proposition}
\begin{proof}
张量商映射 \(A\otimes A\to(A/I)\otimes(A/I)\) 的核为
\(I\otimes A+A\otimes I\)，故第一条件使余乘法下降；其余两条件分别使余单位与对极下降。所有公理在商中继承。对每个 \(R\)，商到 \(R\) 的代数同态就是那些零化 \(I\) 的 \(G(R)\) 元素，以上三式保证它们对乘法、单位和逆封闭。
\end{proof}
例如 \(k[t]\) 上的 \(\Delta(t)=t\otimes1+1\otimes t\) 描述加法群 \(\mathbb G_a\)。特征 \(p\) 中 \(I=(t^p)\) 满足上述条件，因为二项式中间系数为零；所得闭子群概形有非零幂零坐标，不能仅由其 \(k\) 点辨认。

\subsection{一般线性群及可对角化群}
\label{b5:subsec:2102:03}

一般线性群的坐标代数为
\[
k[\operatorname{GL}_n]
=k[x_{ij},d]/\bigl(d\det(x_{ij})-1\bigr).
\]
对任意 \(R\)，它的 \(R\) 点就是可逆矩阵。矩阵乘法给出
\[
\Delta(x_{ij})=\sum_\ell x_{i\ell}\otimes x_{\ell j},
\qquad \varepsilon(x_{ij})=\delta_{ij},
\]
对极则是逆矩阵的条目，由伴随矩阵与行列式逆表示。公式 \(\det(AB)=\det A\det B\) 保证 \(d\) 的余乘法为 \(d\otimes d\)。关系 \(\det X=1\) 又给闭子群 \(\operatorname{SL}_n\)。
\begin{definition}\label{b5:def:diagonalizable-group}
设 \(k\) 为域，\(M\) 为有限生成交换群。群代数 \(k[M]\) 以符号 \(e^m\) 为基，乘法为 \(e^me^n=e^{m+n}\)，余乘法为 \(\Delta(e^m)=e^m\otimes e^m\)。相应的交换仿射群概形记为 \(D(M)\)，称为分裂\term[keduijiaohua-qun]{可对角化群}[diagonalizable group]。当 \(M=\mathbb Z^r\) 时，它是分裂\term[daishu-huanmian]{代数环面}[algebraic torus] \(\mathbb G_m^r\)。
\end{definition}
事实上 \(D(M)(R)=\operatorname{Hom}_{\mathrm{grp}}(M,R^\times)\)，因为代数同态必须把 \(e^m\) 送到相容的单位。特别地，
\[
k[\mathbb G_m]=k[t,t^{-1}],\qquad
D(\mathbb Z/n)=\mu_n,\quad k[\mu_n]=k[t]/(t^n-1).
\]
若 \(k\) 特征为 \(p\)，则 \(t^p-1=(t-1)^p\)，\(\mu_p\) 的坐标环不约化；即使 \(k\) 代数闭，它的 \(k\) 点也只有单位。与此相对，常数群 \(C_p\) 的坐标环为 \(k^{C_p}\)，约化且有 \(p\) 个 \(k\) 点。这是两个不同群概形，不能只因为都带下标 \(p\) 就认同。

\ProseParagraphBreak
可逆矩阵的射影类组成\term[sheying-yiban-xianxingqun]{射影一般线性群}[projective general linear group]
\(\operatorname{PGL}_n\)。作为概形，它是射影矩阵空间
\(\mathbb P_k^{n^2-1}\) 中行列式不为零的仿射开子概形
\(D_+(\det)\)，乘法由矩阵乘法给出，逆由伴随矩阵给出。
在任意交换 \(k\) 代数 \(R\) 上，射影类可在开覆盖上用可逆矩阵表示，重叠处的代表相差一个单位标量；当 \(R\) 为域时，每个类都有全局矩阵代表。
下面的局部计算说明，特殊线性群到这个群的商映射也有坐标环层面的含义。
\begin{proposition}\label{b5:prop:special-projective-quotient}
设 \(k\) 为域，\(n\ge2\)。取矩阵的射影类给出群概形态射
\[
q:\operatorname{SL}_n\longrightarrow\operatorname{PGL}_n.
\]
该态射有限且忠实平坦，核为中心子群概形
\(\mu_n\)，其元素以标量矩阵嵌入。对任意有限维 \(k\) 空间 \(V\)，一个群概形态射
\(r:\operatorname{SL}_n\to\operatorname{GL}(V)\)
唯一地因子化为某个群概形态射
\(\bar r:\operatorname{PGL}_n\to\operatorname{GL}(V)\) 与 \(q\) 的复合，当且仅当 \(r\) 在 \(\mu_n\) 上平凡。
因此 \(\operatorname{PGL}_n\) 是此中心商，记为
\(\operatorname{SL}_n/\mu_n\)。
\end{proposition}
\begin{proof}
在 \(\operatorname{PGL}_n\) 的一个标准仿射开集
\(U=\operatorname{Spec}B\) 上，将某个非零射影坐标归一化为一，得到通用矩阵 \(M\)，其行列式 \(D\) 在 \(B\) 中可逆。
\(q^{-1}(U)\) 上的矩阵恰为 \(tM\)，条件是 \(t^nD=1\)，故其坐标环为
\[
C=B[t]/(t^n-D^{-1}).
\]
它是以 \(1,t,\ldots,t^{n-1}\) 为基的自由 \(B\) 模，而且 \(t\) 可逆。
这些开集覆盖目标，所以 \(q\) 有限且忠实平坦；特别地，它在概形上满射。
单位射影类的原像是 \(tI\)、\(t^n=1\)，即 \(\mu_n\)。
在上述开集上，\(C\otimes_BC\) 中的两个标量参数 \(t,u\) 满足同一个 \(n\) 次幂方程，故 \(z=u/t\) 满足 \(z^n=1\)。这给出
\[
C\otimes_BC\cong C[z]/(z^n-1),\qquad u=tz.
\]
这些同构在重叠处相容，因而
\(\operatorname{SL}_n\times\mu_n\cong
\operatorname{SL}_n\times_{\operatorname{PGL}_n}\operatorname{SL}_n\)，
其中两投影分别为 \((g,z)\mapsto g\) 和 \((g,z)\mapsto gz\)。

还须验证因子化保持正则性，而不只是在点群上取商。
在上述坐标环中，\(\mu_n\) 的作用对应
\[
C\longrightarrow C\otimes_k k[z]/(z^n-1),\qquad
t\longmapsto t\otimes z,\quad b\longmapsto b\otimes1.
\]
若 \(c=\sum_{i=0}^{n-1}b_it^i\) 在此作用下不变，比较
\(1,z,\ldots,z^{n-1}\) 的系数便知 \(b_i=0\) 对所有 \(i>0\) 成立。
因此不变子环恰为 \(B\)，这一计算不要求 \(n\) 在 \(k\) 中可逆。
当 \(r\) 在 \(\mu_n\) 上平凡时，\(r\) 的矩阵条目及行列式逆在该作用下不变，因而都是 \(B\) 中的函数。
它们在 \(U\) 上定义到 \(\operatorname{GL}(V)\) 的态射。
重叠处的两组函数拉回后相同，而这些自由扩张是忠实平坦的，故原函数也相同；于是局部态射拼成唯一的 \(\bar r\)。
在 \(q\times q\) 下拉回后，\(\bar r\) 的乘法相容性就是 \(r\) 的乘法相容性，忠实平坦性再使该等式在目标上成立。
故 \(\bar r\) 是群概形态射。反向若已有因子化，则 \(\ker q=\mu_n\) 显然作用平凡。
\end{proof}
